\subsection{Normalization with a denominator floor}
\label{sec:rmsnorm}

We use the explicit convention
\begin{equation}
 \mathcal N_{\epsilon}(a)_i
 =\frac{a_i}{\sqrt{\epsilon+n^{-1}\sum_{j=1}^n a_j^2}},
 \qquad |a_j|\le R.
 \label{eq:rmsdefinition}
\end{equation}
The original RMSNorm definition has $\epsilon=0$ in this formula
\cite[Eq.~(4)]{zhang2019rmsnorm}. The authors' released implementation
instead adds its numerical stabilizer outside the square root
\cite{zhang2019rmscode}. Llama's reference implementation uses the
inside-root convention in \eqref{eq:rmsdefinition}\cite{llamarmscode}.
These are different functions. Every theorem below refers to
\eqref{eq:rmsdefinition}. A gain of absolute value at most $G$ simply
multiplies the output radius by $G$.

\begin{theorem}[A normalization factory with a denominator floor]
\label{thm:rmsfloor}
Suppose that $R>0$, $\epsilon\ge0$, and $\lambda>0$ are known rational
numbers, and that
\[
 v=\epsilon+\frac1n\sum_{j=1}^n a_j^2\ge\lambda.
\]
Replace $\lambda$ by $\max\{\lambda,\epsilon\}$, and put
\[
 \kappa=\max\{1,R^2/\lambda\},
 \qquad R_{\rm out}=\frac{2R}{\sqrt\lambda}.
\]
Assume that a request at coordinate $j$ returns an independent sign of
mean $a_j/R$. There is an exact sign sampler of mean
$\mathcal N_\epsilon(a)_i/R_{\rm out}$ whose expected number of such
requests is at most
\begin{equation}
 1+228\kappa^2.
 \label{eq:rmschildbound}
\end{equation}
If $B_N$ includes the binary lengths of the rational parameters and
$\lceil\log_2(n+2)\rceil$, its expected auxiliary bit work is
$O(\kappa^2B_N^4)$. The constant in this bound is absolute. No numerical
evaluation of the vector $a$ is required.
\end{theorem}

\paragraph{Proof idea.}
Two independent signs give a disagreement coin whose bias records the
average squared coordinate. A linear factory moves this bias away from its
endpoint, and a negative-binomial generating function supplies the reciprocal
square root. Multiplying by a fresh numerator sign completes the normalized
coordinate.

\begin{proof}
An independent product of two signs with mean $a_j/R$ has mean
$(a_j/R)^2$. It does not give a coin of that probability. The latter
function vanishes when the source bias is zero, an interior point of
the sign-coin parameter interval. A finite Bernoulli factory for that
square coin would therefore be impossible. We instead use the
disagreement coin
\[
 q=\frac12\left(1-\frac1n\sum_j(a_j/R)^2\right).
\]
It costs two source requests at one uniform coordinate. This exact
coin is available even when all coordinates vanish.

Set
\begin{equation}
 U=\epsilon+R^2,\qquad
 d=\frac{\lambda}{2U},\qquad
 C=\frac{2R^2}{U-\lambda/2},\qquad
 \zeta=\frac{\lambda}{2U-\lambda}.
 \label{eq:rmsparameters}
\end{equation}
The promise implies $0<\lambda\le U$. In particular $0<d\le1/2$,
$C\le4$, and
\begin{equation}
 r=Cq\le1-\zeta,
 \qquad \frac C\zeta=\frac{4R^2}{\lambda}.
 \label{eq:rmsamplifier}
\end{equation}
If $C\le1$, obtain an $r$-coin by thinning the $q$-coin. If $C>1$,
use the linear Bernoulli factory with supplied slack
$\min\{\zeta,1/2\}$.

We use the following published input-count result: for $Cp\le1-e$,
Huber's linear factory is exact and uses at most $9.5C/e$ source
coins in expectation \cite[Theorem~1]{huber2016linear}. Its
finite-bit implementation in the parameter range used here is
recorded below. Thus an $r$-coin costs at most $76\kappa$ $q$-coins
in expectation. This bound also covers thinning.

Let
\[
 c_k=\binom{2k}{k}4^{-k},\qquad
 \Pr(K=k)=\sqrt d\,c_k(1-d)^k,\quad k\ge0.
\]
Independently draw a known coin of probability $1/\sqrt2$. If it
fails, return $A=0$. Otherwise return $A=1$ exactly when $K$
independent $r$-coins are all one. For $K=0$, the empty conjunction
is one. The binomial generating function gives
\begin{align}
 \Pr(A=1)
 &=\frac1{\sqrt2}
   \sqrt{\frac d{1-(1-d)r}}
   =\frac{\sqrt\lambda}{2\sqrt v},
 \label{eq:rmsattenuation}\\
 \mathbb E K
 &=\frac{1-d}{2d}
   =\frac U\lambda-\frac12
   \le\frac32\kappa.
 \label{eq:rmsarity}
\end{align}
The algebra in the first line uses
$(1-d)C=2R^2/U$ and
$1-(1-d)Cq=v/U$.

On $A=1$, return a fresh sign of mean $a_i/R$. On $A=0$, return
a fair sign. Its mean is
\[
 \frac{a_i}{R}\frac{\sqrt\lambda}{2\sqrt v}
 =\frac{\mathcal N_\epsilon(a)_i}{R_{\rm out}}.
\]
The fair branch completes the output law to a sign on every run.

The number of $r$-coins is at most $K$. Each $q$-coin uses two
child signs. The uniform conditional input bound for the amplifier
therefore gives
\[
 1+2(76\kappa)(3\kappa/2)=1+228\kappa^2.
\]
This argument charges a fresh factory invocation before its own
randomness is exposed. It does not require the cost of that
invocation to be independent of its output. All component
procedures terminate almost surely, and $K$ is finite almost
surely. This proves exactness, termination, and
\eqref{eq:rmschildbound}.

The arithmetic bound is proved below.
\end{proof}

For recursive composition, the possibly irrational radius
$R_{\rm out}$ can be replaced by a rational one without changing
the child bound. Choose a power of two $\overline R$ with
$R_{\rm out}\le\overline R<2R_{\rm out}$. Its exponent is
found by comparisons of the rational quantities
$\overline R^2\lambda$ and $4R^2$. With probability
$R_{\rm out}/\overline R$, invoke the normalization sampler;
otherwise return a fair sign. The resulting mean is
$\mathcal N_\epsilon(a)_i/\overline R$. This known algebraic
coin has a certified square-root implementation using
$O(B_N^4)$ expected bit work and introduces no child requests.
All later attention and row samplers may therefore use rational
radius bounds.

\subsubsection{Finite-bit implementation}

\begin{lemma}[Half-order negative-binomial sampling]
\label{lem:rmshalfnb}
For rational $0<d\le1/2$, the law
$\Pr(K=k)=\sqrt d\,c_k(1-d)^k$ can be sampled exactly with
$O(d^{-1}(B+\log_2(d^{-1}+2))^4)$ expected bit work, where $B$
bounds the input length of $d$.
\end{lemma}
\paragraph{Proof idea.}
A geometric proposal acquires the required central-binomial factor
through a product of rational acceptance coins. Most unsuccessful products
stop early. Their square-root mean length compensates for the square-root
acceptance probability.

\begin{proof}
Propose $G\ge0$ with law $d(1-d)^k$. Accept it with probability
\[
 c_G=\prod_{j=1}^G\frac{2j-1}{2j},
\]
stopping this product of rational coins at its first failure.
The mean acceptance probability is $\sqrt d$. Thus the accepted
law is the claimed law. The elementary bound
$c_j\le1/\sqrt{j+1}$ follows by induction from
$c_{j+1}/c_j=(2j+1)/(2j+2)$. Conditional on $G=k$, the expected
number of product tests is at most
\[
 \sum_{j=0}^{k-1}c_j\le2\sqrt k.
\]
An unconditional trial uses at most $2/\sqrt d$ tests in
expectation. Dividing by its success probability bounds the
total expected number of tests by $2/d$.

Geometric inversion generates $G$ from
\[
 G=\left\lfloor\frac{-\log V}{-\log(1-d)}\right\rfloor,
 \qquad V\sim\operatorname{Unif}(0,1).
\]
The proof of Lemma~\ref{lem:productaritybits} performs this
inversion with certified logarithms. There the rate is
$-\log(1-p_j)$; here it is $-\log(1-d)\in[d,\log2]$. The
leading-zero count of $V$ has a geometric tail, and the bound $6e$ there
on the chance of lying within $e$ of a positive integer needs only
$e\le1/2$ and a rate times $e$ at most one. Another
$\lceil\log_2(d^{-1}+2)\rceil$ working bits absorb the division by
the rate. Consequently extra refinement has a geometric tail, and a
certified $m$-bit logarithm costs a cubic number of bit operations
in its input and output lengths. Rational product coins use the
usual uniform-bit comparison.

Both the proposal law and the accepted law have geometric
tails on scale $d^{-1}$. Their logarithmic moments, together
with the same bound for the stopped product tests, show that
the expected cost per test is absorbed by the fourth power in
the stated bound. This reasoning counts uniform bits,
counters, rational arithmetic, and precision refinement.
\end{proof}

\begin{lemma}[Finite-bit linear factories with arbitrary known multiplier]
\label{lem:general-linear-bits}
\label{lem:rmslinearbits}
Let $C>1$ and $0<e\le1/2$ be known rationals, and suppose $Cp\le1-e$.
Huber's linear factory has an implementation with
$O(Ce^{-1}B^4)$ expected auxiliary bit work, where $B\ge16$ includes
the parameter lengths, source-address length, and
$\log_2(1/(C-1)+2)$. Source invocations are counted separately.
\end{lemma}
\paragraph{Proof idea.}
Huber's counter has a harmonic function equal to the desired output
bias. Its global source bound controls the early stages, while the rapidly
decreasing probability of reaching a late stage pays for that stage's larger
counters and finer arithmetic.

\begin{proof}
Use Huber's stage algorithm \cite[Theorem~1]{huber2016linear}.
At stage $j$, let $e_j=2^{-j}e$,
$C_j=C\prod_{r<j}(1+e_r/2)<2C$, and $k_j=(23/5)/e_j$.
Initially the integer counter is $i=1$. While $0<i<k_j$, draw a source bit $X$ and
an independent positive geometric variable $G$ of parameter
$(C_j-1)/C_j$, and replace $i$ by $i-1+(1-X)G$. At $i=0$, return
one. At an upper exit, continue with probability $(1+e_j/2)^{-i}$;
on continuation replace $C_j$ by $C_j(1+e_j/2)$ and $e_j$ by $e_j/2$.
If continuation fails, return zero. The source bound for this
algorithm is at most $9.5C/e$.

Its exactness also follows from the counter identity. At a within-stage
step the conditional expectation of $(C_jp)^i$ is unchanged, because
$\mathbb E(C_jp)^G=(C_j-1)p/(1-p)$. At an upper exit the continuation
probability times $(C_{j+1}p)^i$ equals $(C_jp)^i$. The slack promise
is preserved. The finite stopping bounds below imply almost-sure
termination, at which the harmonic function is the returned bit;
bounded convergence therefore gives bias $Cp$.
At an upper exit the probability of continuing is at
most $\exp(-46/25)<1/6$, so the probability of reaching stage $j$ is
at most $6^{-j}$. The stopped-drift calculation, including the geometric
overshoot, bounds the expected update count from any entry state by
\[
 \frac{(C_j-1)\lceil k_j\rceil+C_j}{e_j}
 \le14C e_j^{-2}.
\]
If the entry counter already exceeds the threshold, no updates are
needed before the exit test. A positive geometric variable is sampled
by rational-parameter inversion, so its numerical magnitude is not
charged as a unary loop. All arithmetic and certified logarithm comparisons use dyadic
fixed-point intervals with outward rounding and guarded summation.
The required numerical estimates hold uniformly over the stages. A geometric draw has parameter at least $(C-1)/(2C)$, so
its logarithm has moments $O((B+j)^r)$ for each fixed $r$. At an
upper exit, memorylessness bounds the overshoot by a fresh geometric
excess. Conditioning on continuation multiplies its mass by a
decreasing exponential, which improves this tail. This remains true
at subsequent stage entries.

At requested precision $m$, the product defining $C_j$ is recomputed
with $O(B+j+m)$ working bits and $j$ guarded multiplications.
Subtracting one needs only $O(B)$ additional guard bits, since
$C_j-1\ge C-1$. Geometric inversion uses the logarithm procedure
in Lemma~\ref{lem:rmshalfnb}. For a continuation coin, if
$i e_j>3(m+1)$ its probability is below $2^{-m}$; otherwise its
exponent has $O(B+j+m)$ relevant bits. Binary range reduction and
certified logarithm and exponential series cost a cubic number of
operations in those lengths. Fresh uniform comparison bits give a
geometric refinement tail. These bounds give stage work
$O(Ce_j^{-2}(B+j)^4)$. To sharpen the total, let $N_j$ count updates
at stage $j$. Each update requests exactly one source coin, so Huber's
global source bound gives $\sum_j\mathbb E N_j\le9.5C/e$. At the
start of an update its counter is below $k_j$, hence has length
$O(B+j)$. Its fresh geometric draw and subsequent comparison have
conditional expected auxiliary work $O((B+j)^4)$ uniformly over the
past. Predictable charging bounds the update work by
$O(\sum_j(B+j)^4\mathbb E N_j)$.

Choose $J=\lceil\log(1/e)/\log(3/2)\rceil=O(B)$. The stages
$j\le J$ have total work $O(Ce^{-1}B^4)$ by the global source bound.
For the remaining stages, the reach and stopped-drift estimates give
\[
 \sum_{j>J}(B+j)^4\mathbb E N_j
 \le14Ce^{-2}\sum_{j>J}(2/3)^j(B+j)^4
 =O(Ce^{-1}B^4).
\]
Exit tests have one invocation per reached stage. The conditional
overshoot log-moment estimate and reach probability $6^{-j}$ give a
further $O(B^4)$ term, which is absorbed by the claimed bound.
The construction and its exactness are those of Huber's factory; the
calculation supplies the bit bound at arbitrary known multiplier.
\end{proof}

For the floor parameters, $C/\zeta=4R^2/\lambda\le4\kappa$.
With supplied slack $e=\min\{\zeta,1/2\}$ and $C\le4$, one has
$C/e=O(\kappa)$. Lemma~\ref{lem:general-linear-bits} therefore costs
$O(\kappa B_N^4)$ auxiliary work per amplified coin. There are at
most $3\kappa/2$ such coins in expectation. Predictable charging
bounds their local work by $O(\kappa^2B_N^4)$; source invocations
remain counted separately. Lemma~\ref{lem:rmshalfnb} costs
$O(\kappa B_N^4)$ for the negative-binomial draw. Known gates,
coordinate choices, and the output sign contribute $O(B_N^4)$.
This proves the bit bound in Theorem~\ref{thm:rmsfloor}.
